% Generated draft: requires mathematical and source review.
\documentclass[11pt]{article}
\usepackage[margin=25mm]{geometry}
\usepackage{fontspec}
\setmainfont{Latin Modern Roman}
\usepackage{amsmath,amssymb,mathtools,unicode-math}
\setmathfont{Latin Modern Math}
\usepackage{enumitem,needspace,xurl,hyperref}
\setlength{\parindent}{0pt}
\setlength{\parskip}{5pt}
\setlength{\emergencystretch}{3em}
\newcommand{\N}{\mathbb{N}}
\newcommand{\Z}{\mathbb{Z}}
\newcommand{\Q}{\mathbb{Q}}
\newcommand{\R}{\mathbb{R}}
\newcommand{\C}{\mathbb{C}}
\newcommand{\F}{\mathbb{F}}
\newcommand{\A}{\mathbb{A}}
\newcommand{\PP}{\mathbb{P}}
\newcommand{\E}{\mathbb{E}}
\newcommand{\eps}{\varepsilon}
\newcommand{\dd}{\,\mathrm{d}}
\begin{document}
\section*{UnsolvedMath ID: 700456}
\textbf{Source problem code:} catalogue entry
\subsection{Definitions and mathematical statement}
A loop is a nonempty set \(L\) with a binary multiplication \(\cdot\), an element \(e\), and binary division operations \(\backslash\) and \(/\), such that, for all \(x,y\in L\),
\[
e\cdot x=x=x\cdot e,\qquad x\cdot(x\backslash y)=y,\qquad x\backslash(x\cdot y)=y,\qquad (y/x)\cdot x=y,\qquad (y\cdot x)/x=y.
\]
For this draft, interpret nilpotence as central nilpotence. The center \(Z(L)\) consists of the elements \(z\in L\) such that, for all \(x,y\in L\),
\[
\begin{gathered}
z\cdot x=x\cdot z,\\
(z\cdot x)\cdot y=z\cdot(x\cdot y),\\
(x\cdot z)\cdot y=x\cdot(z\cdot y),\\
(x\cdot y)\cdot z=x\cdot(y\cdot z).
\end{gathered}
\]
Define the upper central series by \(Z_0(L)=\{e\}\) and by taking \(Z_{i+1}(L)\) to be the inverse image of \(Z(L/Z_i(L))\) under the quotient map. The loop is centrally nilpotent if \(Z_c(L)=L\) for some integer \(c\geq 0\).

In the signature \((\cdot,\backslash,/,e)\), let \(\operatorname{Id}(L)\) be the set of all identities \(s=t\) between terms that hold under every assignment of their variables in \(L\). A finite basis for the laws of \(L\), relative to the loop axioms, is a finite subset \(B\subseteq\operatorname{Id}(L)\) such that every loop satisfying all identities in \(B\) satisfies every identity in \(\operatorname{Id}(L)\).

\textbf{Requested construction.} Construct a finite centrally nilpotent loop \(L\) whose laws have no finite basis. Explicitly, require that for every finite subset \(B\subseteq\operatorname{Id}(L)\), there exist a loop \(M\) and an identity \(s=t\in\operatorname{Id}(L)\) such that \(M\) satisfies every identity in \(B\), but \(s=t\) fails under some assignment in \(M\). The witnessing loop \(M\) is not required to be finite.

\textbf{Unresolved source qualifications.} The supplied statement says only ``finite nilpotent loop'' and ``basis for its laws.'' It does not specify the notion of loop nilpotence, the signature of the identities, or the ambient class for equational consequence. The definitions above make explicit the central-nilpotence and full-loop-signature interpretation used in this draft; confirmation of these conventions against the intended source is required.
\subsection{Short English statement}
Construct a finite nilpotent loop whose identities cannot all follow from any finite collection of its identities. This draft interprets nilpotence as central nilpotence and uses multiplication, both division operations, and the identity constant; these conventions require source confirmation.
\subsection{Sources}
\begin{itemize}
\item Supplied local UnsolvedMath registry record: permanent ID \(700456\), exact source problem code \texttt{catalogue entry}, title \emph{Kourovka Notebook Problem 8.4}.
\item Supplied attribution: M. R. Vaughan-Lee, proposed in 1982; original problem number 8.4 in \emph{The Kourovka Notebook (Unsolved Problems in Group Theory)}, 8th issue (1982).
\item The supplied background also identifies the 21st issue edition (2026), \texttt{arXiv:1401.0300v48}. The supplied collection description names E. I. Khukhro and V. D. Mazurov as editors. No external source was independently consulted.
\item The registry labels the problem \texttt{open}; this is supplied curation, not an independently verified assessment of its mathematical or literature status. No separate literature references, literature-source links, or literature-check date are supplied.
\end{itemize}
\end{document}
